% fig/diels_alder_dpo.tex
% Input-only TikZ file.
% Main file should load:
% \usepackage{tikz}
% \usepackage{xparse}
% \usepackage{xcolor}
% \usetikzlibrary{positioning,backgrounds,fit}

\begin{tikzpicture}[
    bondstyle/.style={line width=1pt},
    breakbond/.style={red!75!black,line width=2pt},
    formbond/.style={green!55!black,line width=2pt},
    interfacebond/.style={gray!65,line width=3pt},
    atomnode/.style={
        draw,
        fill=white,
        inner sep=0.5pt,
        circle,
        minimum width=14pt
    },
    interfaceatom/.style={
        draw,
        fill=gray!12,
        inner sep=0.5pt,
        circle,
        minimum width=14pt
    },
    maparrow/.style={->,>=stealth,line width=0.9pt},
    boxstyle/.style={
        draw,
        gray,
        rounded corners,
        inner sep=0.55cm,
        rectangle
    }
]

% ------------------------------------------------------------
% Local parameters
% ------------------------------------------------------------
\def\td{1.2}
\def\graphwidth{4.2}
\def\graphheight{2.5}
\def\colspace{2.0}
\def\rowspace{3.0}

\def\alphaleft{0}
\def\alphacenter{0}
\def\alpharight{0}
\def\alphahost{0}
\def\alphacontext{0}
\def\alphaproduct{0}

% ------------------------------------------------------------
% Helper: double bond
% ------------------------------------------------------------
\NewDocumentCommand{\insertdoublebond}{ O{line width=1pt} O{line width=1pt} m m}{
    \draw[#1] (#3) --
        ++(#4+90:0.05) --
        ++(#4:\td);
    \draw[#2] (#3) --
        ++(#4-90:0.05) --
        ++(#4:\td);
}

% ------------------------------------------------------------
% Coordinates for the graph layout
% Reactive core: 1,2,3,4,6,7
% Context atoms in the application example: 5 and 8
% ------------------------------------------------------------
\newcommand{\insertgraphcoords}[3]{
    \path (#1) coordinate(atom5#2) --
        ++(#3+30:\td) coordinate(atom2#2) --
        ++(#3+90:\td) coordinate(atom3#2) --
        ++(#3+30:\td) coordinate(atom4#2) --
        ++(#3-30:\td) coordinate(atom7#2) --
        ++(#3-90:\td) coordinate(atom6#2) --
        ++(#3-150:\td) coordinate(atom1#2);
    \path (atom7#2) -- ++(#3+30:\td) coordinate(atom8#2);
}

% ------------------------------------------------------------
% L: left-hand side rule graph
% ------------------------------------------------------------
\newcommand{\insertleftgraph}[1]{
    \insertgraphcoords{#1}{L}{\alphaleft}

    \draw[bondstyle] (atom2L) -- (atom3L);
    \insertdoublebond[breakbond][bondstyle]{atom1L}{\alphaleft+150}
    \insertdoublebond[breakbond][bondstyle]{atom3L}{\alphaleft+30}
    \insertdoublebond[breakbond][bondstyle]{atom6L}{\alphaleft+90}

    \foreach \n in {1,2,3,4,6,7}{
        \node[atomnode] at (atom\n L) {\n};
    }
}

% ------------------------------------------------------------
% K: interface graph
% ------------------------------------------------------------
\newcommand{\insertcentergraph}[1]{
    \insertgraphcoords{#1}{K}{\alphacenter}

    % \draw[interfacebond] (atom1K) -- (atom6K);
    \draw[interfacebond] (atom6K) -- (atom7K);
    % \draw[interfacebond] (atom7K) -- (atom4K);
    \draw[interfacebond] (atom4K) -- (atom3K);
    \draw[interfacebond] (atom3K) -- (atom2K);
    \draw[interfacebond] (atom2K) -- (atom1K);

    \foreach \n in {1,2,3,4,6,7}{
        \node[interfaceatom] at (atom\n K) {\n};
    }
}

% ------------------------------------------------------------
% R: right-hand side rule graph
% Green bonds highlighted at 1-6 and 4-7
% ------------------------------------------------------------
\newcommand{\insertrightgraph}[1]{
    \insertgraphcoords{#1}{R}{\alpharight}

    \draw[bondstyle] (atom2R) -- (atom1R);
    \draw[formbond]  (atom1R) -- (atom6R);
    \draw[bondstyle] (atom6R) -- (atom7R);
    \draw[formbond]  (atom7R) -- (atom4R);
    \draw[bondstyle] (atom4R) -- (atom3R);

    \insertdoublebond[formbond][bondstyle]{atom2R}{\alpharight+90}

    \foreach \n in {1,2,3,4,6,7}{
        \node[atomnode] at (atom\n R) {\n};
    }
}

% ------------------------------------------------------------
% G: host graph
% ------------------------------------------------------------
\newcommand{\inserthostgraph}[1]{
    \insertgraphcoords{#1}{G}{\alphahost}

    \draw[bondstyle] (atom5G) -- (atom2G) -- (atom3G);
    \draw[bondstyle] (atom7G) -- (atom8G);

    \insertdoublebond[breakbond][bondstyle]{atom1G}{\alphahost+150}
    \insertdoublebond[breakbond][bondstyle]{atom3G}{\alphahost+30}
    \insertdoublebond[breakbond][bondstyle]{atom6G}{\alphahost+90}

    \foreach \n in {1,...,8}{
        \node[atomnode] at (atom\n G) {\n};
    }
}

% ------------------------------------------------------------
% D: context graph
% ------------------------------------------------------------
\newcommand{\insertcontextgraph}[1]{
    \insertgraphcoords{#1}{D}{\alphacontext}

    % \draw[interfacebond] (atom1D) -- (atom6D);
    \draw[interfacebond] (atom6D) -- (atom7D);
    % \draw[interfacebond] (atom7D) -- (atom4D);
    \draw[interfacebond] (atom4D) -- (atom3D);
    \draw[interfacebond] (atom3D) -- (atom2D);
    \draw[interfacebond] (atom2D) -- (atom1D);

    \draw[bondstyle] (atom5D) -- (atom2D);
    \draw[bondstyle] (atom7D) -- (atom8D);

    \foreach \n in {1,...,8}{
        \node[interfaceatom] at (atom\n D) {\n};
    }
}

% ------------------------------------------------------------
% H: product graph
% Green bonds highlighted at 1-6 and 4-7
% ------------------------------------------------------------
\newcommand{\insertproductgraph}[1]{
    \insertgraphcoords{#1}{H}{\alphaproduct}

    \draw[bondstyle] (atom5H) -- (atom2H);
    \draw[bondstyle] (atom7H) -- (atom8H);

    \draw[bondstyle] (atom2H) -- (atom1H);
    \draw[formbond]  (atom1H) -- (atom6H);
    \draw[bondstyle] (atom6H) -- (atom7H);
    \draw[formbond]  (atom7H) -- (atom4H);
    \draw[bondstyle] (atom4H) -- (atom3H);

    \insertdoublebond[formbond][bondstyle]{atom2H}{\alphaproduct+90}

    \foreach \n in {1,...,8}{
        \node[atomnode] at (atom\n H) {\n};
    }
}

% ------------------------------------------------------------
% Box helpers
% ------------------------------------------------------------
\newcommand{\insertboxtop}[3]{
    \node[boxstyle,
        fit=(atom5#1) (atom8#1) (atom1#1) (atom4#1)
    ] (box#2) {};
    \node[fill=white,above=-0.18cm of box#2] {\strut #3};
}

\newcommand{\insertboxbottom}[3]{
    \node[boxstyle,
        fit=(atom5#1) (atom8#1) (atom1#1) (atom4#1)
    ] (box#2) {};
    \node[fill=white,below=-0.18cm of box#2] {\strut #3};
}

% ------------------------------------------------------------
% Draw graphs
% ------------------------------------------------------------

% Upper row: rule L <- K -> R
\insertleftgraph{0,0}
\insertcentergraph{\graphwidth+\colspace,0}
\insertrightgraph{2*\graphwidth+2*\colspace,0}

% Lower row: application G <- D -> H
\inserthostgraph{0,-\rowspace-\graphheight}
\insertcontextgraph{\graphwidth+\colspace,-\rowspace-\graphheight}
\insertproductgraph{2*\graphwidth+2*\colspace,-\rowspace-\graphheight}

% Boxes
\begin{scope}[on background layer]
    \insertboxtop{L}{L}{L}
    \insertboxtop{K}{K}{K}
    \insertboxtop{R}{R}{R}

    \insertboxbottom{G}{G}{G}
    \insertboxbottom{D}{D}{D}
    \insertboxbottom{H}{H}{H}
\end{scope}

% ------------------------------------------------------------
% Rule arrows: L <- K -> R
% Now from box edge to box edge
% ------------------------------------------------------------
\draw[maparrow,shorten <=4pt,shorten >=4pt]
    (boxK.west) -- node[above] {$l$} (boxL.east);

\draw[maparrow,shorten <=4pt,shorten >=4pt]
    (boxK.east) -- node[above] {$r$} (boxR.west);

% ------------------------------------------------------------
% Application arrows: G <- D -> H
% ------------------------------------------------------------
\draw[maparrow,shorten <=4pt,shorten >=4pt]
    (boxD.west) -- node[above] {$g$} (boxG.east);

\draw[maparrow,shorten <=4pt,shorten >=4pt]
    (boxD.east) -- node[above] {$h$} (boxH.west);

% ------------------------------------------------------------
% Vertical maps m, m', m''
% Box-to-box only
% ------------------------------------------------------------
\draw[maparrow,shorten <=4pt,shorten >=4pt]
    (boxL.south) -- node[right] {$m$} (boxG.north);

\draw[maparrow,shorten <=4pt,shorten >=4pt]
    (boxK.south) -- node[right] {$m'$} (boxD.north);

\draw[maparrow,shorten <=4pt,shorten >=4pt]
    (boxR.south) -- node[right] {$m''$} (boxH.north);

\end{tikzpicture}